\section{Introduction}
\label{sec:intro}
% numbers finalized from P1 v4 + P4A/E + P4C outputs (2026-10-10).

Web agents have crossed from demos to deployment: they shop, book, and
compare on our behalf, turning the Web's product pages from artifacts that
\emph{inform humans} into signals that \emph{direct machine decisions}. At
the same time, the supply side of those pages is being rewritten. An
estimated 27.5\% of tokens passing common quality filters on the open Web
are now AI-generated, and the share is climbing monthly \cite{russell2026much};
measurement studies document machine-generated catalogs flooding music
platforms \cite{wu2026empirical} and scientific publishing \cite{oh2026science}.
This content is not adversarial in the classical sense: much of it is produced
by ordinary commercial incentives --- search-optimized content farms,
counterfeit-style listing clones, review inflation --- and a whole research
program now studies how to \emph{optimize} content for generative engines
\cite{wu2025what,bagga2025egeo,nimase2026geobench} and how to detect the
result \cite{chu2026geoflag,zheng2026countergeobench,li2026when}. The
inevitable consequence is that agentic systems increasingly browse a Web
whose supply of plausible, keyword-rich, well-rated listings no longer
correlates with actually satisfying the user's goal.

This paper asks a question that, to our knowledge, no prior work has measured
in a controlled setting: \textbf{does ambient, non-targeted synthetic content
silently erode the decision quality of web agents?} Existing evaluations of
agents under Web pollution attack them deliberately: programmatic mini-internets
with adversarially curated falsehoods \cite{shah2026synthetic}, attacker-written
endorsement pages steering search agents \cite{chen2026much}, or prompt
injection and memory poisoning \cite{cao2026ward,yang2026zombie,syros2026muzzle}.
Upstream of the decision, retrieval-level studies show that cumulative AI
pollution degrades search quality \cite{yu2026retrieval}. But ambient
contamination is a different regime: no adversary knows the user's goal, so
no page is crafted to defeat a specific decision. The threat is statistical,
not surgical --- and whether it materially matters for agent \emph{economics}
(choice quality, wasted spend, calibrated confidence) is unknown.

Measuring this regime requires solving a confound that has kept prior work at
the retrieval layer: if we simply compare agent success before and after
contaminating a benchmark, any degradation could stem from task difficulty
rather than choice quality. We therefore build \textbf{SynthFlood}, a paired,
dose-controlled protocol on top of WebShop's official artifacts
\cite{yao2022webshop}. Each task renders a fixed candidate list in which the
gold product is always present and never replaced; only distractor slots are
progressively replaced by synthetic listings, drawn under a fixed seed so
that lower contamination levels are nested within higher ones. Synthetic
listings come from a three-family generator that mirrors real content-farm
modes --- search-optimized pages, specification clones with altered attributes
and lower prices, and review-inflated plain pages --- with a cross-model
plausibility floor, all produced by open-weight 7--8B models. Agents act in a
ReAct loop over rendered listings, and every purchase is scored by WebShop's
\emph{official} reward function, so regret is measured in the benchmark's own
currency rather than a proxy.

SynthFlood yields three measurements. First, a \emph{forced-choice
susceptibility probe}: one gold product against one synthetic listing, 300
randomized duels, yielding a single interpretable number $\sigma$ --- the
probability that an agent buys the machine-optimized page over the correct
product when nothing else competes. Second, a \emph{dose--response curve}:
success, regret, and calibration as functions of the contamination share
$\varepsilon\in\{0,0.1,0.25,0.5,0.75\}$, with McNemar tests on paired tasks.
Our silent-erosion hypothesis predicts a characteristic signature here:
regret rises steeply while coarse success appears stable, because partial
mismatches (wrong attributes, overpayment) are invisible in success counts.
Third, a \emph{defense battery} evaluated at matched contamination, whose
critical control is a random-defense arm that drops the same number of
candidates uniformly --- separating genuine filtering from mere candidate
removal --- alongside a provenance-only defense that re-ranks by a noisy
verifiable-provenance bit without ever seeing labels.

%FILL: contributions + headline numbers after gates.
Our contributions are:
\begin{itemize}
\item \textbf{A threat model and protocol} for ambient (non-targeted)
      synthetic contamination of agent decision channels, with paired
      dose--response control on official benchmark artifacts.
\item \textbf{A forced-choice probe with channel decomposition} isolating
      \emph{why} machine-optimized listings win: specification clones sit at
      chance ($\sigma_{\text{clone}}=0.523$) while gloss and review
      inflation are rejected ($\sigma_{\text{geo}}=0.351$,
      $\sigma_{\text{reviews}}=0.287$); but price parity alone drops the
      clone win-rate to $0.386$ (paired $p=0.043$), stripping inflated
      ratings alone also yields $0.386$, and removing both drops it to
      $0.250$ --- an additive two-channel decomposition showing the attack
      exploits price and social proof, not perception. A pricing
      dose--response yields a threshold law: any discount $\geq 10\%$
      flips the majority to the synthetic clone ($0.386\!\to\!0.591$,
      $p=0.003$), with no further gain at deeper discounts.
\item \textbf{The first decision-level dose--response evidence} that ambient
      contamination erodes agent economics: at $\varepsilon=0.75$, success
      drops by $23$\,pp ($p=1.3\times10^{-4}$; $p=1.2\times10^{-8}$ pooled
      over two seeds) and regret rises $69\%$ (paired $\Delta=+0.181$,
      95\% CI $[+0.095,+0.264]$), while confidence rises and calibration
      error doubles --- errors that no success metric would surface.
\item \textbf{A defense analysis with an ordering principle}: an
      agent-visible noisy provenance badge eliminates the regret excess,
      is unchanged from 15\% to 65\% badge forgery, and a text-free
      structural detector statistically matches a perfect-provenance
      oracle ($p=0.21$); but consuming the same bit as a soft ranking
      bonus fails until its weight exceeds the rating scale
      ($\lambda\leq 4$ all fail, $\lambda=8$ recovers fully) --- and a
      naive popularity ranking \emph{backfires}, raising synthetic
      purchases to 67.5\%.
\end{itemize}
We release all code, contamination pools, and full trajectory logs. The
remainder of the paper formalizes the threat model
(\S\ref{sec:problem}), presents the protocol (\S\ref{sec:method}), reports
results (\S\ref{sec:results}), discusses implications and limitations
(\S\ref{sec:discussion}), and concludes (\S\ref{sec:conclusion}).

\paragraph{What is brave and new.}
This paper proposes that the Web's next decade will be shaped not by
adversarial attacks on agents but by \emph{ambient} commercial synthetic
content --- and that the resulting damage is invisible to the metrics we use
to evaluate agents. What is new: (i) the first decision-level dose--response
evidence that ambient contamination erodes agent economics on official
benchmark artifacts; (ii) a channel-decomposed forced-choice probe that
locates the erosion mechanism in the agent's \emph{preference function}
(price and social proof), not its perception, with a pricing threshold law
showing a 10\% discount is sufficient to flip the majority; (iii) a
diagnostic triplet --- flat success, rising regret, rising confidence ---
that operators can monitor; and (iv) the demonstration that a single noisy
provenance bit recovers the full loss when consumed as an \emph{ordering
constraint} --- but not as a soft score --- where content-based detection
cannot. What is brave: we argue provenance-ordered ranking should become a
first-class interface contract between the Web and the agents acting on it,
and we give the community the protocol to falsify or refine that claim.
